% 38_body.tex — 산출물 ㊳ 시리즈A IR 1페이저 · 재무모델 3케이스 (빈 템플릿 / 완성 예시 공용 본문) — gen_product_templates.py 가 Product_specs.py 에서 생성
% 드라이버: 38_IR1페이저재무모델_빈템플릿.tex (\wbblanktrue) / 38_IR1페이저재무모델_완성예시.tex (\wbblankfalse — 워크북 §X.5 확정 후)
\documentclass[10pt,a4paper]{article}
\usepackage[a4paper,margin=16mm,top=14mm,bottom=20mm]{geometry}
\def\DocNum{38}\def\DocDay{8}\def\DocIdx{3}
\def\DocTitle{시리즈A IR 1페이저 · 재무모델 3케이스}
\def\DocSub{Series A One-pager \& 3-Case Financial Model}
\def\DocVariant{\B{빈 템플릿}{딥게이지 완성 예시}}
\usepackage{wbstyle}
\def\DocCirc{㊳}   % newunicodechar(㉛~㊵ 폴백) 뒤에 정의해야 활성 문자로 읽힌다
\begin{document}
\docheader
 {\Bg{[회사명]}{딥게이지 · GAUGE}}
 {\Bg{[v1.x / 2028-05]}{v1.2 / 2028-05-12 (D+802)}}
 {\Bg{[작성자 — CEO]}{강민수 CEO\rd{7mm}}}
 {\Bg{[검토자 — 기존 투자자]}{서영채(노들투자파트너스)\rd{7mm}}}

%% ------------------------------------------------------------------ 1
\secbar{1. IR 1페이저}
\guide{문제 / 제품 / 트랙션 / 유닛이코노믹스 / 팀 / 조달과 사용처 — 각 칸 3줄. **숫자 표기 규칙**: 계약 기준과 검산값을 병기하거나 각주에 분모를 쓴다.}
{\footnotesize
\begin{tabularx}{\linewidth}{|L{26mm}|Y|}\hline
\hd{요소} & \hd{내용} \\ \hline
\ifwbblank
\rd{11mm}문제(한 문장 + 숫자) &  \\ \hline
\rd{11mm}제품(무엇이 달라지는가) &  \\ \hline
\rd{11mm}트랙션(고객·ARR·리텐션) &  \\ \hline
\rd{11mm}유닛이코노믹스(GM·CAC·burn multiple) &  \\ \hline
\rd{11mm}팀 &  \\ \hline
\rd{11mm}조달 규모·사용처·마일스톤 &  \\ \hline
\else
문제(한 문장 + 숫자) & 검사원은 검사보다 검사기록에 시간을 쓴다 — 작성 47분(전기 31) / 통보 3.7일 / 8D 11.5시간 / 클레임 2,700만(G-33~36). 초기 SAM 2,150 · 확장 11,800 \\ \hline
제품(무엇이 달라지는가) & GAUGE — 사진·음성으로 검사기록을 만들고 판정 초안과 근거를 남긴다. 판정을 대신하지 않는다(자동판정 기본 OFF, 채택률 26\%는 설계값). Flow 애드온으로 8D·포털 제출 \\ \hline
트랙션(고객·ARR·리텐션) & 유료 58개사 · ARR 15.68억(검산 — 계약 기준 18.9, 온담 라이선스 3.78만 ARR / 온담 제외 11.9) · GRR 87 / NDR 104\%(D+435 21사 기준) · 27Q1 코호트 M12 71\% \\ \hline
유닛이코노믹스(GM·CAC·burn multiple) & GM 46.9\%(계약 MRR 기준 54\%) — AI 원가 44.0\% · CAC blended 740만 → payback 10.7개월(계약 기준 8.0) · burn multiple 1.20 · 순소진 8,400만 · 현금 7.52억(8.9개월) \\ \hline
팀 & 27명(개발 14 / 세일즈 6 / CS 4 / 경영지원 3) · CEO 품질 현장 · CTO 2-pass · CPO 47건 3분류 · 온담 투입 상한 6 FTE \\ \hline
조달 규모·사용처·마일스톤 & 80억 / pre 320억 · 개발·eval 32 / 세일즈·CS 20 / 온프렘·보안 8 / 온담 적자·제품화 5.82 / 운전자본 14.18 · 12개월 후 ARR 21.8억(검산) · 온담 4세트 측정 완료 · NDR 106\% \\ \hline
\fi
\end{tabularx}}

%% ------------------------------------------------------------------ 2
\secbar{2. 트랙션 표 — 두 값 병기}
\guide{유료 58개사 · ARR 11.9억(온담 제외) / 18.9억(IR, 온담 포함) / 15.68억(검산) · NDR 104\% · GM 54/46.9\% · burn multiple 1.20.}
{\footnotesize
\begin{tabularx}{\linewidth}{|L{40mm}|Y|Y|Y|}\hline
\hd{지표} & \hd{IR 기재값} & \hd{검산값} & \hd{분모·각주} \\ \hline
\ifwbblank
\rd{9mm}ARR &  &  &  \\ \hline
\rd{9mm}GM &  &  &  \\ \hline
\rd{9mm}NDR / GRR &  &  &  \\ \hline
\rd{9mm}CAC payback &  &  &  \\ \hline
\rd{9mm}burn multiple &  &  &  \\ \hline
\rd{9mm}집중 리스크 &  &  &  \\ \hline
\else
ARR & 18.9억(온담 포함) & 15.68억(라이선스만) · 온담 제외 11.9 · 순SaaS 10.3 & 온담 연 청구 7.0억 중 ARR 성격은 3.78억(라인 70 × 45만 × 12) \\ \hline
GM & 54\% & 46.9\% & 분모 계약 MRR 9,917 vs 순SaaS 8,583 · COGS 4,560(AI 3,780) \\ \hline
NDR / GRR & 104\% / 87\% & 동일 & D+435 유료 21사 MRR 2,150만 · 이탈 −280 · 확장 +366 \\ \hline
CAC payback & 8.0개월 & 10.7개월 & CAC 740 ÷ (ACV/12 × GM) — 2,052·54\% vs 1,776·46.9\% \\ \hline
burn multiple & 1.20 & 1.20(순SaaS) · 계약 기준 0.99 & 8,400 × 3 ÷ 분기 순증 2.10억 — 덱이 순SaaS를 쓴 유일한 지표 \\ \hline
집중 리스크 & 37.0\% & 24.1\% · 개발 용량 82.1\% & 7.0/18.9 · 3.78/15.68 · 138/168 \\ \hline
\fi
\end{tabularx}}

%% ------------------------------------------------------------------ 3
\secbar{3. 재무모델 3케이스 — Base / Bull / Bear}
\guide{24개월. 온담 딜 포함·제외를 분리하고 매출 인식 시점(라이선스·구축·운영)을 명시. 각 케이스의 가정 3개.}
{\footnotesize
\begin{tabularx}{\linewidth}{|L{40mm}|Y|Y|Y|}\hline
\hd{항목} & \hd{Bear} & \hd{Base} & \hd{Bull} \\ \hline
\ifwbblank
\rd{9mm}가정 1 (고객 수) &  &  &  \\ \hline
\rd{9mm}가정 2 (ACV·NDR) &  &  &  \\ \hline
\rd{9mm}가정 3 (온담 딜) &  &  &  \\ \hline
\rd{9mm}ARR 12개월 후 &  &  &  \\ \hline
\rd{9mm}ARR 24개월 후 &  &  &  \\ \hline
\rd{9mm}월 순소진 / 런웨이 &  &  &  \\ \hline
\rd{9mm}필요 조달 &  &  &  \\ \hline
\else
가정 1 (고객 수) & 58 → 70 → 82 & 58 → 84 → 118 & 58 → 96 → 150 \\ \hline
가정 2 (ACV·NDR) & 2,052 유지 · 100\% & 2,052 → 2,150 → 2,250 · 104 → 106\% & 2,052 → 2,300 → 2,500 · 106 → 110\% \\ \hline
가정 3 (온담 딜) & 거절 → 대한기공 2.4억(코너스톤 소멸) & 수용(3분류) — 라이선스 3.78 ARR · 구축 3.0/6.12 1년차 · 운영 2.22 통과 & 수용 + 식품 3사(24개월 차 +2.0억) \\ \hline
ARR 12개월 후 & 16.8억 & 21.8억(IR 표기 25.0) & 25.9억 \\ \hline
ARR 24개월 후 & 19.2억 & 30.3억(IR 표기 33.5) & 43.3억 \\ \hline
월 순소진 / 런웨이 & 3.4억 / 8.4개월(현금 28.4) & 2.8억 / 13.2개월(현금 36.8) & 1.9억 / 25.1개월(현금 47.6) \\ \hline
필요 조달 & 시리즈B 앞당김(20개월) 또는 채용 동결 · 브릿지 30억 & 시리즈B 150억 착수 18~20개월(ARR 26억) & 선택적 — 24개월 후 \\ \hline
\fi
\end{tabularx}}

%% ------------------------------------------------------------------ 4
\secbar{4. 사용처와 마일스톤}
\guide{80억을 어디에(개발·세일즈·온프렘·eval 확장) · 24개월 마일스톤 · 다음 라운드 전제.}
{\footnotesize
\begin{tabularx}{\linewidth}{|L{40mm}|L{22mm}|Y|L{22mm}|}\hline
\hd{사용처} & \hd{금액} & \hd{마일스톤(무엇이 되면)} & \hd{시점} \\ \hline
\ifwbblank
\rd{9mm} &  &  &  \\ \hline
\rd{9mm} &  &  &  \\ \hline
\rd{9mm} &  &  &  \\ \hline
\rd{9mm} &  &  &  \\ \hline
\rd{9mm} &  &  &  \\ \hline
\rd{9mm} &  &  &  \\ \hline
\else
개발·eval 확장(14 → 21 → 40 · DS-v2 · 온담 4세트 80건 측정 1,200만 포함) & 32.00 & 온담 4세트 측정 완료 · FP ≤ 15\% 라인 과반 & D+1000 + 8주 / 2029-Q2 \\ \hline
세일즈·CS(6 → 14 · 4 → 9) & 20.00 & 유료 84사 · NDR 106\% & 2029-05 \\ \hline
온프렘·보안(ISMS-P · 218 조건부 62) & 8.00 & ISMS-P 획득 · 온프렘 2사 & 2029-Q4 \\ \hline
온담 구축 적자 3.12 · 제품화 27 mm 2.70 & 5.82 & 커버리지 61.4 → 95\% · 검증 6종 통과 & 2029-Q1 \\ \hline
운전자본·예비 & 14.18 & 런웨이 12개월 유지 & 상시 \\ \hline
계 & 80.00 &  &  \\ \hline
\fi
\end{tabularx}}

%% ------------------------------------------------------------------ 5
\secbar{5. 리스크 공개}
\guide{집중 리스크(3표기) · AI 원가율 · 채택률 · 도메인 확장 · 핵심 인력. 숨기지 않은 리스크가 IC에서 방어된다.}
{\footnotesize
\begin{tabularx}{\linewidth}{|L{40mm}|Y|Y|}\hline
\hd{리스크} & \hd{현재 상태(숫자)} & \hd{완화} \\ \hline
\ifwbblank
\rd{10mm} &  &  \\ \hline
\rd{10mm} &  &  \\ \hline
\rd{10mm} &  &  \\ \hline
\rd{10mm} &  &  \\ \hline
\rd{10mm} &  &  \\ \hline
\else
집중 리스크 & 37.0 / 24.1 / 82.1\% & 3분류(86 mm) · 상한 6 FTE · SOW 7.08억 · 대한기공 병행 → 24.1 → 18\% \\ \hline
AI 원가율 & 44.0\%(순SaaS) · 단가 67원 & 조건부 2-pass · 단가 45원 → 32\% \\ \hline
채택률·FP & 26\%(㉮) · FP 21.3\% & DS-v2 오후 세트 · FP 예산 ≤ 15\% 라인만 ON \\ \hline
도메인 확장(식품) & 온담 4세트 미측정 & 측정 계획 80건 · 8주 · 1,200만 — 인수 기준 조항(㊵ 항목 4) \\ \hline
핵심 인력 & 개발 14명 중 시니어 4명 & 온담 투입 상한 6 FTE(45\%) — 자동차 릴리스 분기 1회 보장 \\ \hline
\fi
\end{tabularx}}

%% ------------------------------------------------------------------ 6
\secbar{6. IC 예상 반대신문 8문과 답}
\guide{'집중 리스크가 37\%인데 24.1\%라고요?' '채택률 22\%면 제품이 안 쓰이는 것 아닌가요?' '온담을 거절하면 코너스톤은?' 등 8문. 답에 분모를 쓴다.}
{\footnotesize
\begin{tabularx}{\linewidth}{|L{8mm}|Y|Y|}\hline
\hd{\#} & \hd{질문} & \hd{답(분모 포함)} \\ \hline
\ifwbblank
\rd{10mm} &  &  \\ \hline
\rd{10mm} &  &  \\ \hline
\rd{10mm} &  &  \\ \hline
\rd{10mm} &  &  \\ \hline
\rd{10mm} &  &  \\ \hline
\rd{10mm} &  &  \\ \hline
\rd{10mm} &  &  \\ \hline
\rd{10mm} &  &  \\ \hline
\else
1 & 집중 리스크가 37\%인데 24.1\%라고요? & 분모가 다르다 — 18.9는 온담 연 청구 7억 전액, 15.68은 라이선스 3.78만. 세 번째 숫자 82.1\%(138/168)를 우리가 먼저 말한다 \\ \hline
2 & 채택률 26\%면 제품이 안 쓰이는 것 아닌가요? & 분모 ㉮ AI 노출 건. 자동판정 기본 OFF 설계 — 커널. 지표는 NS-1 주간 승인 검사기록(58사 27,600건/주). FP 21.3\%가 원인, DS-v2가 조치 \\ \hline
3 & 온담을 거절하면 코너스톤은? & 소멸(선행 조건). 한강 단독 → 옵션풀 협상력 감소. BATNA 대한기공 2.4억 × 3년 \\ \hline
4 & GM 54\%인데 AI 원가율 44\%면? & 54\%는 계약 MRR 9,917 분모, 44\%는 순SaaS 8,583. 같은 분모면 46.9\% / 38.1\%. 덱에 병기 \\ \hline
5 & 사고 이후 재발했습니까? & 0건 — 감시 대시보드 위의 0건('무신호' 아님). 재발 방지 6건 증빙(㉞ 항목 8) = 제218문항 답변 \\ \hline
6 & 구축비 3억이 6.12억이면 적자 아닌가요? & 1년차 −4.04억, 3년 공헌 5.52억(SOW 7.08억 별도). 그래서 배상 상한 10\%·커스텀 유상 \\ \hline
7 & 개발 14명 중 몇 명이 온담에? & 분류 후 86 mm ÷ 12 = 7.2 FTE(51\%) — 상한 6 FTE로 계약. 전부 수용이면 14명 9.9개월 \\ \hline
8 & NDR 104\%의 분모는? & D+435 유료 21사 MRR 2,150만. 이탈 −280(GRR 87) · 확장 +366(라인 증설 + Flow) \\ \hline
\fi
\end{tabularx}}

\end{document}